\section{学习工程的边界、规则与证据}

\topic{学习对象不是代码量}
操作系统连接处理器、内存和设备，同时向程序提供受保护的抽象。学习项目真正
要获得的是机制之间的因果关系：复位向量位于固定地址，页表同时承担翻译和
权限，中断改变执行流，调度器维护的也远不止一个函数指针数组。代码只是把
这些关系变成可执行证据。

\begin{keyidea}
每次只增加一个新的硬件或软件不变量。先让 ROM 可执行，再让串口可观察，再
加载下一阶段；不要同时引入模式切换、ELF、页表和 C++ 运行时。
\end{keyidea}

\topic{项目边界}
本项目允许 QEMU TCG 模拟硬件，因而宿主机不需要是 x86-64。QEMU 不提供
启动软件；传给 \code{-bios} 的 128 KiB 镜像完全由项目生成。宿主 Python
负责调度 CMake、CTest、NASM、LLD 和 QEMU，但不会进入目标镜像。

\begin{center}
\begin{tikzpicture}[os diagram]
  \node[os block] (source) {NASM / C++20\\项目源代码};
  \node[os state,right=of source] (image) {ROM / 磁盘\\可审计镜像};
  \node[os hardware,right=of image] (qemu) {QEMU TCG\\硬件模型};
  \node[os state,below=of qemu] (evidence) {串口、状态、\\测试证据};
  \draw[os arrow] (source) -- (image);
  \draw[os arrow] (image) -- (qemu);
  \draw[os arrow] (qemu) -- (evidence);
\end{tikzpicture}
\end{center}

\topic{软件工程闭环}
需求定义“必须做到什么”；架构说明模块和交接；ADR 记录关键取舍；代码实现
机制；测试保存可重复证据；发布记录对应一个可复现里程碑。完成阶段必须同时
满足构建、静态审计、宿主测试、整机测试、失败路径和文档六个方面。

\topic{语言与代码规则}
目标代码使用显式宽度类型，如 \code{std::uint32_t} 和
\code{std::uint64_t}。硬件寄存器按协议宽度表示；内部计数、容量和地址若
没有更窄的硬件约束，优先使用 64 位。类成员访问显式写 \code{this->}；
常量采用 \code{OS_模块_功能} 的全大写命名；头文件与实现分离；中文注释
解释约束和原因，不复述语法。

\topic{操作系统所处的位置}
处理器只能取指、译码和执行，内存只保存位模式，设备也只按照各自的寄存器
协议改变状态。它们不会主动提供“进程”“文件”“终端”或“线程”这些抽象。
操作系统位于裸硬件与应用程序之间，把有限而异构的硬件资源组织成稳定的
软件接口，同时负责隔离互不信任的执行实体。

\begin{pdfdiagram}[node distance=7mm]
  \node[os state,minimum width=55mm] (app) {应用程序\\进程、文件、套接字};
  \node[os block,below=of app,minimum width=55mm] (kernel)
    {内核机制\\地址空间、中断、调度、系统调用};
  \node[os block,below=of kernel,minimum width=55mm] (driver)
    {设备驱动\\寄存器协议与并发状态};
  \node[os hardware,below=of driver,minimum width=55mm] (hardware)
    {处理器、内存控制器、定时器、磁盘、串口};
  \draw[os arrow] (app) -- node[right]{受控接口} (kernel);
  \draw[os arrow] (kernel) -- node[right]{统一资源模型} (driver);
  \draw[os arrow] (driver) -- node[right]{总线事务} (hardware);
\end{pdfdiagram}
\diagramnote{抽象由内核状态、硬件状态和失败处理共同维持。}

内核是操作系统中始终拥有最高硬件权限的部分，但完整的操作系统还包括启动
固件、引导代码、运行库、系统服务、命令解释器和维护工具。本项目从复位开始，
是因为只有亲自建立这些前置条件，才能分清哪些能力来自处理器，哪些来自平台，
哪些才是内核提供的抽象。

\topic{操作系统形成的历史脉络}
早期计算机一次只运行一个程序。操作员手工装入程序、设置开关并取走结果，
昂贵的处理器却常常等待人和机械外设。批处理系统把一组作业连续装入，监控
程序负责在作业之间交接机器，由此出现了最初的常驻控制软件。

处理器速度继续增长以后，磁带、打印机和磁盘的延迟成为主要浪费。多道程序
设计让一个作业等待 I/O 时切换到另一个作业；中断提供异步完成通知，内存
保护防止作业互相破坏。分时系统进一步利用周期性时钟中断快速轮换任务，使
多个终端用户获得近似独占机器的交互体验。虚拟内存随后把进程看到的地址与
实际物理位置分离，分页硬件既扩大了可用地址空间，也成为隔离和共享的基础。

这些机制并非按教科书目录一次设计完成。它们分别回应了设备缓慢、资源昂贵、
程序错误、多人共享和软件规模增长等现实压力。本书沿着依赖关系重新组织它们：
先取得执行权与可观测性，再建立内存、异常和并发，最后形成进程、文件和用户态。

\topic{启动链路中的责任边界}\par
\begin{osgridlongtable}{L{0.20\textwidth}L{0.32\textwidth}L{0.34\textwidth}}
\textbf{层次} & \textbf{承担的责任} & \textbf{尚未提供的能力}\\ \hline
\endfirsthead
\textbf{层次} & \textbf{承担的责任} & \textbf{尚未提供的能力}\\ \hline
\endhead
平台模型 & 地址译码、RAM、ROM、设备和复位行为 & 固件策略、磁盘格式、内核接口\\ \hline
固件 & 从复位环境建立最小执行环境 & 通用进程、文件系统、用户态 ABI\\ \hline
引导阶段 & 定位并校验内核，准备交接数据 & 长期设备管理与应用隔离\\ \hline
内核 & 内存、异常、中断、调度、驱动与系统调用 & 应用业务逻辑\\ \hline
用户空间 & 服务、命令和应用 & 任意访问物理内存或设备寄存器\\ \hline
\end{osgridlongtable}

明确边界可以防止“模拟器替我们完成了启动”的错觉。QEMU 只实现选定硬件的
行为；ROM 中每个字节、每次模式切换、每张页表和每个磁盘读取协议都由项目
代码建立。

\topic{机制、策略与可验证证据}
机制描述系统能够执行的动作，例如切换页表、屏蔽中断或把任务放入就绪队列；
策略决定何时以及对谁执行这些动作，例如页面替换算法和调度优先级。两者分离
以后，可以先证明底层动作正确，再在不破坏硬件契约的前提下更换策略。

一次成功启动只是示例，不是证明。工程证据至少包含可复现构建、二进制布局
审计、成功路径、故障注入、超时边界以及输出判定。代码、文档和测试共同描述
同一个系统：代码给出机制，文档给出约束与理由，测试验证外部可观察行为。
